\section{Introduction to linear transformations}
\textbf{A matrix as a rule.} An $m\times n$ matrix acts on an
$n$-coordinate input and produces an $m$-coordinate output. The
\emph{domain} is $\R^n$, the \emph{codomain} is $\R^m$, and the
\emph{range} is the subset of outputs actually attained. A linear
transformation satisfies
\[
 T(cu+dv)=cT(u)+dT(v)
\]
for all inputs $u,v$ and scalars $c,d$. In particular $T(0)=0$.
To find inputs producing $b$, solve $Ax=b$. To find inputs sent to
zero, solve $Ax=0$. These are different questions: existence concerns
whether $b$ is attained; uniqueness concerns whether it is attained
by only one input.

\practice{1.8.P1}{Read matrix dimensions from input and output}
Five input coordinates require five columns, so that each column
receives one coefficient. Two output coordinates require two rows,
one row computation per output entry. Thus $A$ is $2\times5$.

\practice{1.8.P2}{Interpret the sign change}
Multiplication by $\begin{bmatrix}1&0\\0&-1\end{bmatrix}$ takes
$(x_1,x_2)$ to $(x_1,-x_2)$. Horizontal position is unchanged and
vertical position is reversed, so this is reflection across the $x_1$-axis.
Points on that axis are fixed.

\practice{1.8.P3}{Map a segment starting at zero}
A point of the segment has the form $x=tu$ with $0\le t\le1$.
Linearity gives $T(x)=tT(u)$, which has exactly the form of a
point on the segment from $0$ to $T(u)$ with the same parameter.
Conversely each such output comes from $tu$. Thus the image is the
whole segment, possibly collapsed to one point if $T(u)=0$.

\exercise{1.8.1}{Multiply by a scalar matrix}
Here $A=2I_2$. Hence
\[
 T(1,-3)^T=(2,-6)^T,\qquad T(a,b)^T=(2a,2b)^T.
\]
Each output entry is twice the corresponding input entry. Geometrically
this doubles lengths from the origin without changing direction.

\exercise{1.8.2}{A uniform contraction}
Here $A=\tfrac13 I_3$, so multiplying each input coordinate by
$1/3$ gives
\[
 T(3,6,-9)^T=(1,2,-3)^T,\qquad
 T(a,b,c)^T=(a/3,b/3,c/3)^T.
\]
The factor is $1/3$, not $3$: it is the value on each diagonal entry.


\exercise{1.8.3}{Solve for an input with the specified image}

\goal Solve $Ax=b$ and decide whether the input is unique. Put the target vector in the last column and apply the same row operations to the whole augmented matrix.

\[ [A\mid b]=\begin{bmatrix}1 & 0 & -3 & -2\\-3 & 1 & 6 & 3\\2 & -2 & -1 & -1\end{bmatrix}.\]
Use $R_{2}\leftarrow R_{2}-(-3)R_{1},\quad R_{3}\leftarrow R_{3}-(2)R_{1}$. Then
\[ \begin{bmatrix}1 & 0 & -3 & -2\\0 & 1 & -3 & -3\\0 & -2 & 5 & 3\end{bmatrix}.\]
Use $R_{3}\leftarrow R_{3}-(-2)R_{2}$. Then
\[ \begin{bmatrix}1 & 0 & -3 & -2\\0 & 1 & -3 & -3\\0 & 0 & -1 & -3\end{bmatrix}.\]
Use $R_{3}\leftarrow -1R_{3},\quad R_{1}\leftarrow R_{1}-(-3)R_{3},\quad R_{2}\leftarrow R_{2}-(-3)R_{3}$. Then
\[ \begin{bmatrix}1 & 0 & 0 & 7\\0 & 1 & 0 & 6\\0 & 0 & 1 & 3\end{bmatrix}.\]

The reduced equations give $x=(7,6,3)^T$. Every one of the three variable columns has a pivot, so the solution is unique.

As an independent check, multiplying the displayed particular input by the original $A$ gives $\begin{bmatrix}-2\\3\\-1\end{bmatrix}$, exactly the required target.


\exercise{1.8.4}{Solve for an input with the specified image}

\goal Solve $Ax=b$ and decide whether the input is unique. Put the target vector in the last column and apply the same row operations to the whole augmented matrix.

\[ [A\mid b]=\begin{bmatrix}1 & -2 & 3 & -6\\0 & 1 & -3 & -4\\2 & -5 & 6 & -5\end{bmatrix}.\]
Use $R_{3}\leftarrow R_{3}-(2)R_{1}$. Then
\[ \begin{bmatrix}1 & -2 & 3 & -6\\0 & 1 & -3 & -4\\0 & -1 & 0 & 7\end{bmatrix}.\]
Use $R_{1}\leftarrow R_{1}-(-2)R_{2},\quad R_{3}\leftarrow R_{3}-(-1)R_{2}$. Then
\[ \begin{bmatrix}1 & 0 & -3 & -14\\0 & 1 & -3 & -4\\0 & 0 & -3 & 3\end{bmatrix}.\]
Use $R_{3}\leftarrow - \frac{1}{3}R_{3},\quad R_{1}\leftarrow R_{1}-(-3)R_{3},\quad R_{2}\leftarrow R_{2}-(-3)R_{3}$. Then
\[ \begin{bmatrix}1 & 0 & 0 & -17\\0 & 1 & 0 & -7\\0 & 0 & 1 & -1\end{bmatrix}.\]

The reduced equations give $x=(-17,-7,-1)^T$. Every variable is a pivot variable, so this input is unique.

As an independent check, multiplying the displayed particular input by the original $A$ gives $\begin{bmatrix}-6\\-4\\-5\end{bmatrix}$, exactly the required target.


\exercise{1.8.5}{Solve for an input with the specified image}

\goal Solve $Ax=b$ and decide whether the input is unique. Put the target vector in the last column and apply the same row operations to the whole augmented matrix.

\[ [A\mid b]=\begin{bmatrix}1 & -5 & -7 & -2\\-3 & 7 & 5 & -2\end{bmatrix}.\]
Use $R_{2}\leftarrow R_{2}-(-3)R_{1}$. Then
\[ \begin{bmatrix}1 & -5 & -7 & -2\\0 & -8 & -16 & -8\end{bmatrix}.\]
Use $R_{2}\leftarrow - \frac{1}{8}R_{2},\quad R_{1}\leftarrow R_{1}-(-5)R_{2}$. Then
\[ \begin{bmatrix}1 & 0 & 3 & 3\\0 & 1 & 2 & 1\end{bmatrix}.\]

The reduced equations give $x=(3-3t,1-2t,t)^T$, $t\in\R$. Here $x_3=t$ is free. For example $t=0$ gives $(3,1,0)^T$. Varying $t$ gives different inputs with the same output, so uniqueness fails.

As an independent check, multiplying the displayed particular input by the original $A$ gives $\begin{bmatrix}-2\\-2\end{bmatrix}$, exactly the required target.


\exercise{1.8.6}{Solve for an input with the specified image}

\goal Solve $Ax=b$ and decide whether the input is unique. Put the target vector in the last column and apply the same row operations to the whole augmented matrix.

\[ [A\mid b]=\begin{bmatrix}1 & -3 & 2 & 1\\3 & -8 & 8 & 6\\0 & 1 & 2 & 3\\1 & 0 & 8 & 10\end{bmatrix}.\]
Use $R_{2}\leftarrow R_{2}-(3)R_{1},\quad R_{4}\leftarrow R_{4}-(1)R_{1}$. Then
\[ \begin{bmatrix}1 & -3 & 2 & 1\\0 & 1 & 2 & 3\\0 & 1 & 2 & 3\\0 & 3 & 6 & 9\end{bmatrix}.\]
Use $R_{1}\leftarrow R_{1}-(-3)R_{2},\quad R_{3}\leftarrow R_{3}-(1)R_{2},\quad R_{4}\leftarrow R_{4}-(3)R_{2}$. Then
\[ \begin{bmatrix}1 & 0 & 8 & 10\\0 & 1 & 2 & 3\\0 & 0 & 0 & 0\\0 & 0 & 0 & 0\end{bmatrix}.\]

The reduced equations give $x=(10-8t,3-2t,t)^T$, $t\in\R$. Here $x_3=t$ is free. Taking $t=0$ gives $(10,3,0)^T$, but other values give additional inputs. The zero rows impose no extra restrictions.

As an independent check, multiplying the displayed particular input by the original $A$ gives $\begin{bmatrix}1\\6\\3\\10\end{bmatrix}$, exactly the required target.


\exercise{1.8.7}{The domain and codomain of a six-by-five matrix}
There are five columns, so $Ax$ requires $x\in\R^5$; there are
six rows, so $Ax\in\R^6$. Therefore $\boxed{a=5,b=6}$.
The range can be smaller than $\R^6$; the question asks for the codomain.

\exercise{1.8.8}{Construct the required shape}
A map $\R^5\to\R^7$ implemented by $Ax$ needs $\boxed{7\text{ rows
and }5\text{ columns}}$. Five entries are consumed and seven entries
are produced. Reversing these numbers would define a different map.


\exercise{1.8.9}{Find every input sent to zero}

\goal Solve $Ax=0$. The right-hand side stays zero under row operations, so we can reduce the coefficient matrix alone.

\[ A=\begin{bmatrix}1 & -3 & 5 & -5\\0 & 1 & -3 & 5\\2 & -4 & 4 & -4\end{bmatrix}.\]
Use $R_{3}\leftarrow R_{3}-(2)R_{1}$. Then
\[ \begin{bmatrix}1 & -3 & 5 & -5\\0 & 1 & -3 & 5\\0 & 2 & -6 & 6\end{bmatrix}.\]
Use $R_{1}\leftarrow R_{1}-(-3)R_{2},\quad R_{3}\leftarrow R_{3}-(2)R_{2}$. Then
\[ \begin{bmatrix}1 & 0 & -4 & 10\\0 & 1 & -3 & 5\\0 & 0 & 0 & -4\end{bmatrix}.\]
Use $R_{3}\leftarrow - \frac{1}{4}R_{3},\quad R_{1}\leftarrow R_{1}-(10)R_{3},\quad R_{2}\leftarrow R_{2}-(5)R_{3}$. Then
\[ \begin{bmatrix}1 & 0 & -4 & 0\\0 & 1 & -3 & 0\\0 & 0 & 0 & 1\end{bmatrix}.\]

The equations are $x_1-4x_3=0$, $x_2-3x_3=0$, $x_4=0$. Let $x_3=t$. Then
\[ x=t\begin{bmatrix}4\\3\\1\\0\end{bmatrix},\qquad t\in\R.\]

Each displayed generator is sent to zero by the original matrix. Conversely, every solution must obey the reduced equations, so no additional inputs are missing.


\exercise{1.8.10}{Find every input sent to zero}

\goal Solve $Ax=0$. The right-hand side stays zero under row operations, so we can reduce the coefficient matrix alone.

\[ A=\begin{bmatrix}3 & 2 & 10 & -6\\1 & 0 & 2 & -4\\0 & 1 & 2 & 3\\1 & 4 & 10 & 8\end{bmatrix}.\]
Use $R_{1}\leftrightarrow R_{2},\quad R_{2}\leftarrow R_{2}-(3)R_{1},\quad R_{4}\leftarrow R_{4}-(1)R_{1}$. Then
\[ \begin{bmatrix}1 & 0 & 2 & -4\\0 & 2 & 4 & 6\\0 & 1 & 2 & 3\\0 & 4 & 8 & 12\end{bmatrix}.\]
Use $R_{2}\leftrightarrow R_{3},\quad R_{3}\leftarrow R_{3}-(2)R_{2},\quad R_{4}\leftarrow R_{4}-(4)R_{2}$. Then
\[ \begin{bmatrix}1 & 0 & 2 & -4\\0 & 1 & 2 & 3\\0 & 0 & 0 & 0\\0 & 0 & 0 & 0\end{bmatrix}.\]

The equations are $x_1+2x_3-4x_4=0$ and $x_2+2x_3+3x_4=0$. Set $x_3=s,x_4=t$. Then
\[ x=s\begin{bmatrix}-2\\-2\\1\\0\end{bmatrix}+t\begin{bmatrix}4\\-3\\0\\1\end{bmatrix},\quad s,t\in\R.\]

Each displayed generator is sent to zero by the original matrix. Conversely, every solution must obey the reduced equations, so no additional inputs are missing.


\exercise{1.8.11}{Test membership in the range}

The target belongs to the range exactly when $Ax=b$ is consistent. Reducing the augmented matrix gives
\[ [A\mid b]\sim \begin{bmatrix}1 & 0 & -4 & 0 & 2\\0 & 1 & -3 & 0 & 1\\0 & 0 & 0 & 1 & 0\end{bmatrix}.\]

The equations give $x_1=2+4t$, $x_2=1+3t$, $x_3=t$, $x_4=0$. In particular $x=(2,1,0,0)^T$ produces $b=(-1,1,0)^T$. Thus the target is in the range. A free variable does not prevent existence; it only prevents uniqueness.


\exercise{1.8.12}{Test membership in the range}

The target belongs to the range exactly when $Ax=b$ is consistent. Reducing the augmented matrix gives
\[ [A\mid b]\sim \begin{bmatrix}1 & 0 & 2 & -4 & 0\\0 & 1 & 2 & 3 & 0\\0 & 0 & 0 & 0 & 1\\0 & 0 & 0 & 0 & 0\end{bmatrix}.\]

The third row says $0=1$, which is impossible. There is no input with image $(-1,3,-1,4)^T$, so the target is not in the range. The fact that it lies in the codomain $\R^4$ is not enough.


\exercise{1.8.13}{Calculate and sketch the transformed vectors}

Multiply each row of $A$ by the input coordinates:
\[ T(u)=\begin{bmatrix}-5\\-2\end{bmatrix},\qquad T(v)=\begin{bmatrix}2\\-4\end{bmatrix}.\]

For a general point, $T(x_1,x_2)=(- x_{1},- x_{2})$. This is reflection through the origin (equivalently a rotation through $180^\circ$). The black arrows show the original vectors; the blue arrows show their images.

\begin{center}\setlength{\unitlength}{1mm}
\begin{picture}(130,90)
\put(5,30){\vector(1,0){118}}
\put(60,5){\vector(0,1){78}}
\put(122,25){$x_1$}\put(62,82){$x_2$}
\put(56,25){\scriptsize$0$}
\put(20,29){\line(0,1){2}}\put(20,25){\makebox(0,0){\scriptsize$-10$}}
\put(40,29){\line(0,1){2}}\put(40,25){\makebox(0,0){\scriptsize$-5$}}
\put(80,29){\line(0,1){2}}\put(80,25){\makebox(0,0){\scriptsize$5$}}
\put(100,29){\line(0,1){2}}\put(100,25){\makebox(0,0){\scriptsize$10$}}
\put(59,10){\line(1,0){2}}\put(55,10){\makebox(0,0){\scriptsize$-5$}}
\put(59,50){\line(1,0){2}}\put(55,50){\makebox(0,0){\scriptsize$5$}}
\put(59,70){\line(1,0){2}}\put(55,70){\makebox(0,0){\scriptsize$10$}}
{\color{black}\linethickness{0.8pt}\put(60,30){\vector(5,2){20}}}
\put(82,40){\makebox(0,0)[l]{\small$u$}}
{\color{black}\linethickness{0.8pt}\put(60,30){\vector(-1,2){8}}}
\put(43,48){\makebox(0,0)[l]{\small$v$}}
{\color{blue}\linethickness{0.8pt}\put(60,30){\vector(-5,-2){20}}}
\put(42,18){\makebox(0,0)[l]{\small$T(u)$}}
{\color{blue}\linethickness{0.8pt}\put(60,30){\vector(1,-2){8}}}
\put(70,16){\makebox(0,0)[l]{\small$T(v)$}}
\end{picture}\end{center}


\exercise{1.8.14}{Calculate and sketch the transformed vectors}

Multiply each row of $A$ by the input coordinates:
\[ T(u)=\begin{bmatrix}10\\4\end{bmatrix},\qquad T(v)=\begin{bmatrix}-4\\8\end{bmatrix}.\]

For a general point, $T(x_1,x_2)=(2 x_{1},2 x_{2})$. This is expansion by a factor of $2$ about the origin. The black arrows show the original vectors; the blue arrows show their images.

\begin{center}\setlength{\unitlength}{1mm}
\begin{picture}(130,90)
\put(5,30){\vector(1,0){118}}
\put(60,5){\vector(0,1){78}}
\put(122,25){$x_1$}\put(62,82){$x_2$}
\put(56,25){\scriptsize$0$}
\put(20,29){\line(0,1){2}}\put(20,25){\makebox(0,0){\scriptsize$-10$}}
\put(40,29){\line(0,1){2}}\put(40,25){\makebox(0,0){\scriptsize$-5$}}
\put(80,29){\line(0,1){2}}\put(80,25){\makebox(0,0){\scriptsize$5$}}
\put(100,29){\line(0,1){2}}\put(100,25){\makebox(0,0){\scriptsize$10$}}
\put(59,10){\line(1,0){2}}\put(55,10){\makebox(0,0){\scriptsize$-5$}}
\put(59,50){\line(1,0){2}}\put(55,50){\makebox(0,0){\scriptsize$5$}}
\put(59,70){\line(1,0){2}}\put(55,70){\makebox(0,0){\scriptsize$10$}}
{\color{black}\linethickness{0.8pt}\put(60,30){\vector(5,2){20}}}
\put(82,40){\makebox(0,0)[l]{\small$u$}}
{\color{black}\linethickness{0.8pt}\put(60,30){\vector(-1,2){8}}}
\put(43,48){\makebox(0,0)[l]{\small$v$}}
{\color{blue}\linethickness{0.8pt}\put(60,30){\vector(5,2){40}}}
\put(102,42){\makebox(0,0)[l]{\small$T(u)$}}
{\color{blue}\linethickness{0.8pt}\put(60,30){\vector(-1,2){16}}}
\put(46,64){\makebox(0,0)[l]{\small$T(v)$}}
\end{picture}\end{center}


\exercise{1.8.15}{Calculate and sketch the transformed vectors}

Multiply each row of $A$ by the input coordinates:
\[ T(u)=\begin{bmatrix}2\\5\end{bmatrix},\qquad T(v)=\begin{bmatrix}4\\-2\end{bmatrix}.\]

For a general point, $T(x_1,x_2)=(x_{2},x_{1})$. This is reflection through the line $x_2=x_1$. The black arrows show the original vectors; the blue arrows show their images.

\begin{center}\setlength{\unitlength}{1mm}
\begin{picture}(130,90)
\put(5,30){\vector(1,0){118}}
\put(60,5){\vector(0,1){78}}
\put(122,25){$x_1$}\put(62,82){$x_2$}
\put(56,25){\scriptsize$0$}
\put(20,29){\line(0,1){2}}\put(20,25){\makebox(0,0){\scriptsize$-10$}}
\put(40,29){\line(0,1){2}}\put(40,25){\makebox(0,0){\scriptsize$-5$}}
\put(80,29){\line(0,1){2}}\put(80,25){\makebox(0,0){\scriptsize$5$}}
\put(100,29){\line(0,1){2}}\put(100,25){\makebox(0,0){\scriptsize$10$}}
\put(59,10){\line(1,0){2}}\put(55,10){\makebox(0,0){\scriptsize$-5$}}
\put(59,50){\line(1,0){2}}\put(55,50){\makebox(0,0){\scriptsize$5$}}
\put(59,70){\line(1,0){2}}\put(55,70){\makebox(0,0){\scriptsize$10$}}
{\color{black}\linethickness{0.8pt}\put(60,30){\vector(5,2){20}}}
\put(82,40){\makebox(0,0)[l]{\small$u$}}
{\color{black}\linethickness{0.8pt}\put(60,30){\vector(-1,2){8}}}
\put(43,48){\makebox(0,0)[l]{\small$v$}}
{\color{blue}\linethickness{0.8pt}\put(60,30){\vector(2,5){8}}}
\put(70,46){\makebox(0,0)[l]{\small$T(u)$}}
{\color{blue}\linethickness{0.8pt}\put(60,30){\vector(2,-1){16}}}
\put(78,24){\makebox(0,0)[l]{\small$T(v)$}}
\end{picture}\end{center}


\exercise{1.8.16}{Calculate and sketch the transformed vectors}

Multiply each row of $A$ by the input coordinates:
\[ T(u)=\begin{bmatrix}0\\4\end{bmatrix},\qquad T(v)=\begin{bmatrix}0\\8\end{bmatrix}.\]

For a general point, $T(x_1,x_2)=(0,2 x_{2})$. This is projection onto the $x_2$-axis followed by multiplication of the remaining coordinate by $2$. The black arrows show the original vectors; the blue arrows show their images.

\begin{center}\setlength{\unitlength}{1mm}
\begin{picture}(130,90)
\put(5,30){\vector(1,0){118}}
\put(60,5){\vector(0,1){78}}
\put(122,25){$x_1$}\put(62,82){$x_2$}
\put(56,25){\scriptsize$0$}
\put(20,29){\line(0,1){2}}\put(20,25){\makebox(0,0){\scriptsize$-10$}}
\put(40,29){\line(0,1){2}}\put(40,25){\makebox(0,0){\scriptsize$-5$}}
\put(80,29){\line(0,1){2}}\put(80,25){\makebox(0,0){\scriptsize$5$}}
\put(100,29){\line(0,1){2}}\put(100,25){\makebox(0,0){\scriptsize$10$}}
\put(59,10){\line(1,0){2}}\put(55,10){\makebox(0,0){\scriptsize$-5$}}
\put(59,50){\line(1,0){2}}\put(55,50){\makebox(0,0){\scriptsize$5$}}
\put(59,70){\line(1,0){2}}\put(55,70){\makebox(0,0){\scriptsize$10$}}
{\color{black}\linethickness{0.8pt}\put(60,30){\vector(5,2){20}}}
\put(82,40){\makebox(0,0)[l]{\small$u$}}
{\color{black}\linethickness{0.8pt}\put(60,30){\vector(-1,2){8}}}
\put(43,48){\makebox(0,0)[l]{\small$v$}}
{\color{blue}\linethickness{0.8pt}\put(60,30){\vector(0,1){16}}}
\put(62,42){\makebox(0,0)[l]{\small$T(u)$}}
{\color{blue}\linethickness{0.8pt}\put(60,30){\vector(0,1){32}}}
\put(62,64){\makebox(0,0)[l]{\small$T(v)$}}
\end{picture}\end{center}


\exercise{1.8.17}{Apply linearity before looking for a matrix}
The given images are $T(u)=(4,1)^T$ and $T(v)=(-1,3)^T$.
Homogeneity and then additivity give
\[
 T(2u)=2T(u)=(8,2)^T,\qquad T(3v)=3T(v)=(-3,9)^T,
\]
\[
 T(2u+3v)=2T(u)+3T(v)=(8,2)^T+(-3,9)^T=(5,11)^T.
\]
We do not need to reconstruct $T$ on every input to answer these questions.

\exercise{1.8.18}{Transfer a graphical linear combination}
In the input diagram, $w=u+2v$: adding two copies of the
left-pointing vector $v$ cancels the horizontal component of $u$
without changing its vertical component. By linearity,
\[
 T(w)=T(u)+2T(v).
\]
In the output diagram, begin at the tip of $T(u)$, append two copies
of $T(v)$ head to tail, and draw the arrow from the origin to the final
point. The blue arrow below gives the result. These are schematic
coordinate drawings of the relation, not new numerical measurements
assigned to the textbook's unlabeled vectors.

\begin{center}\setlength{\unitlength}{1mm}
\begin{picture}(110,75)
\put(5,22){\vector(1,0){98}}
\put(45,5){\vector(0,1){63}}
\put(102,17){$x_1$}\put(47,67){$x_2$}
\put(41,17){\scriptsize$0$}
{\color{black}\linethickness{0.8pt}\put(45,22){\vector(1,1){28}}}
\put(75,52){\makebox(0,0)[l]{\small$u$}}
{\color{black}\linethickness{0.8pt}\put(45,22){\vector(-1,0){14}}}
\put(25,18){\makebox(0,0)[l]{\small$v$}}
{\color{blue}\linethickness{0.8pt}\put(45,22){\vector(0,1){28}}}
\put(39,53){\makebox(0,0)[l]{\small$w$}}
\end{picture}\end{center}

\begin{center}\setlength{\unitlength}{1mm}
\begin{picture}(120,75)
\put(5,25){\vector(1,0){108}}
\put(25,5){\vector(0,1){63}}
\put(112,20){$x_1$}\put(27,67){$x_2$}
\put(21,20){\scriptsize$0$}
{\color{black}\linethickness{0.8pt}\put(25,25){\vector(2,-1){26}}}
\put(53,14){\makebox(0,0)[l]{\small$T(u)$}}
{\color{black}\linethickness{0.8pt}\put(25,25){\vector(16,17){20.8}}}
\put(47.8,49.1){\makebox(0,0)[l]{\small$T(v)$}}
{\color{blue}\linethickness{0.8pt}\put(25,25){\vector(13,6){67.6}}}
\put(94.6,58.2){\makebox(0,0)[l]{\small$T(w)$}}
\end{picture}\end{center}


\exercise{1.8.19}{Use the standard coordinate vectors}
Write $(5,-3)^T=5e_1-3e_2$. Then
\[
 T(5,-3)^T=5(2,5)^T-3(-1,6)^T=(13,7)^T.
\]
More generally, $x=x_1e_1+x_2e_2$ gives
\[
 T(x)=x_1(2,5)^T+x_2(-1,6)^T=(2x_1-x_2,5x_1+6x_2)^T.
\]
The images of the two coordinate vectors determine every image.

\exercise{1.8.20}{Make the two image vectors into columns}
The rule is $T(x)=x_1(-3,5)^T+x_2(7,-2)^T$. Matrix multiplication
forms exactly this column combination, so
\[
 A=\begin{bmatrix}-3&7\\5&-2\end{bmatrix},\qquad T(x)=Ax.
\]
Putting the image vectors into rows instead would give the wrong formula.

\exercise{1.8.21}{Functions, codomains, ranges, and linearity}
\begin{enumerate}[label=(\alph*)]
\item \textbf{True.} A transformation is a function; a linear transformation
is one satisfying additional addition and scaling identities.
\item \textbf{False.} Five columns require inputs in $\R^5$. Three
rows determine the codomain $\R^3$, not the domain.
\item \textbf{False.} The range is $\Col A$, which need not equal
$\R^m$. For the zero matrix the only output is zero.
\item \textbf{True in the setting of this chapter, $\R^n\to\R^m$.}
For any such linear $T$, write $x=\sum_jx_je_j$; then
$T(x)=\sum_jx_jT(e_j)=[T(e_1)\ \cdots\ T(e_n)]x$.
Thus it has a matrix implementation. A transformation on other vector
spaces requires a choice of coordinates before this literal formula applies.
\item \textbf{True.} Additivity and homogeneity imply the displayed
combination rule. Conversely, setting both scalars to $1$ gives
additivity, and setting the second scalar to $0$ gives homogeneity.
Those are exactly the two defining properties.
\end{enumerate}

\exercise{1.8.22}{Separate an existence question from a uniqueness question}
\begin{enumerate}[label=(\alph*)]
\item \textbf{True.} The products $Ax$ are precisely all weighted
combinations of the columns, by the definition of matrix multiplication.
\item \textbf{True.} $A(cu+dv)=cAu+dAv$, by distributivity and scaling
of matrix products.
\item \textbf{False.} Membership of $c$ in the range asks whether
\emph{some} input gives $c$: it is an existence question. A uniqueness
question asks whether two different inputs can give that same output.
\item \textbf{True.} These are the defining properties of a linear map.
\item \textbf{True.} $T(0)=T(0+0)=T(0)+T(0)$; subtracting $T(0)$
from both sides gives $T(0)=0$.
\end{enumerate}

\exercise{1.8.23}{A straight-line graph need not define a linear map}
\step{(a) Put $b=0$.}
For $f(x)=mx$ and any real $c,d,u,v$,
$f(cu+dv)=m(cu+dv)=c(mu)+d(mv)=cf(u)+df(v)$. Hence it is linear.
\step{(b) Inspect zero when $b\ne0$.}
Then $f(0)=b\ne0$, violating the necessary condition that a linear
transformation sends zero to zero.
\step{(c) Explain the terminology.}
In elementary algebra and calculus, ``linear function'' often means a
function whose graph is a straight line, allowing an intercept. Linear
algebra uses the stricter operation-preserving meaning. A nonzero intercept
makes the function affine rather than linear in this stricter sense.

\exercise{1.8.24}{An affine displacement}
For $T(x)=Ax+b$, the zero input gives $T(0)=b$. If $b\ne0$,
this cannot equal zero, whereas every linear transformation sends zero
to zero. Thus the affine map is not linear. The matrix part is linear;
it is the added fixed displacement that violates the condition.

\exercise{1.8.25}{The image of an entire line}
Write the input line as $x=p+tv$ for all $t\in\R$, with $v\ne0$.
Then $T(x)=T(p)+tT(v)$. If $T(v)\ne0$, these outputs form the
line through $T(p)$ in direction $T(v)$. Every point of that line
comes from the same parameter $t$, so the map is onto that line.
If $T(v)=0$, every input on the line maps to the single point $T(p)$.
\answer{The image is a line or a single point.}

\exercise{1.8.26}{Affine combinations describe lines and segments}
\step{(a) Rewrite a line's usual parametrization.}
For distinct $p,q$, the line is $p+t(q-p)$. Distributing gives
$(1-t)p+tq$, for all real $t$. Conversely these formulas are equal
for every $t$, so they describe the same set. If $p=q$, the formula
instead describes the degenerate line consisting of that one point.
\step{(b) Restrict the parameter to the segment.}
For $0\le t\le1$, linearity gives
$T((1-t)p+tq)=(1-t)T(p)+tT(q)$. As $t$ runs over the whole
interval $[0,1]$, the outputs cover the segment between $T(p)$ and
$T(q)$. If these endpoints coincide, that segment is a single point.

\exercise{1.8.27}{A plane can collapse in two ways}
Every input in $P$ is $su+tv$ with $s,t\in\R$. Its image is
$sT(u)+tT(v)$, and conversely every combination of those two image
vectors comes from such an input. Thus
\[
 T(P)=\Span\{T(u),T(v)\}.
\]
If the image vectors are independent, their span is a plane through zero.
If they are dependent but at least one is nonzero, their span is a line
through zero. If both are zero, the span is just $\{0\}$.
\answer{The image is a plane precisely when $T(u),T(v)$ are independent.}

\exercise{1.8.28}{A parallelogram maps into the corresponding parallelogram}
An input in the given parallelogram has the form $au+bv$ with
$0\le a,b\le1$. Its image is $aT(u)+bT(v)$ with the same two
coefficients in the same ranges. This is precisely a point of the
parallelogram generated by $T(u),T(v)$. In fact the image is the
whole of that set, since every allowed pair $a,b$ also gives an input.
The output parallelogram may degenerate to a segment or a point.

\exercise{1.8.29}{Two sketches for reflection across the vertical axis}
Here $T(x_1,x_2)=(-x_1,x_2)$.
For addition, choose $u=(2,1)$, $v=(1,2)$. Then $u+v=(3,3)$,
while $T(u)=(-2,1)$, $T(v)=(-1,2)$, and
$T(u)+T(v)=(-3,3)=T(u+v)$. The paired arrows below have matching
vertical coordinates and opposite horizontal coordinates.

\begin{center}\setlength{\unitlength}{1mm}
\begin{picture}(130,90)
\put(5,15){\vector(1,0){118}}
\put(65,5){\vector(0,1){78}}
\put(122,10){$x_1$}\put(67,82){$x_2$}
\put(61,10){\scriptsize$0$}
{\color{black}\linethickness{0.8pt}\put(65,15){\vector(2,1){24}}}
\put(91,29){\makebox(0,0)[l]{\small$u$}}
{\color{black}\linethickness{0.8pt}\put(65,15){\vector(1,2){12}}}
\put(79,41){\makebox(0,0)[l]{\small$v$}}
{\color{black}\linethickness{0.8pt}\put(65,15){\vector(1,1){36}}}
\put(103,53){\makebox(0,0)[l]{\small$u+v$}}
{\color{blue}\linethickness{0.8pt}\put(65,15){\vector(-2,1){24}}}
\put(26,24){\makebox(0,0)[l]{\small$T(u)$}}
{\color{blue}\linethickness{0.8pt}\put(65,15){\vector(-1,2){12}}}
\put(38,41){\makebox(0,0)[l]{\small$T(v)$}}
{\color{blue}\linethickness{0.8pt}\put(65,15){\vector(-1,1){36}}}
\put(9,54){\makebox(0,0)[l]{\small$T(u+v)$}}
\end{picture}\end{center}

For scaling, take $u=(2,1)$ and $c=2$. Then $T(2u)=(-4,2)=2T(u)$. Both the original ray and its reflected ray double their lengths; their directions are reflected across the vertical axis.

\begin{center}\setlength{\unitlength}{1mm}
\begin{picture}(145,80)
\put(5,15){\vector(1,0){133}}
\put(72,5){\vector(0,1){68}}
\put(137,10){$x_1$}\put(74,72){$x_2$}
\put(68,10){\scriptsize$0$}
{\color{black}\linethickness{0.8pt}\put(72,15){\vector(2,1){24}}}
\put(98,23){\makebox(0,0)[l]{\small$u$}}
{\color{black}\linethickness{0.8pt}\put(72,15){\vector(2,1){48}}}
\put(122,41){\makebox(0,0)[l]{\small$2u$}}
{\color{blue}\linethickness{0.8pt}\put(72,15){\vector(-2,1){24}}}
\put(33,23){\makebox(0,0)[l]{\small$T(u)$}}
{\color{blue}\linethickness{0.8pt}\put(72,15){\vector(-2,1){48}}}
\put(4,43){\makebox(0,0)[l]{\small$2T(u)=T(2u)$}}
\end{picture}\end{center}


\exercise{1.8.30}{A linear map vanishing on a spanning list is zero}
Take any $x\in\R^n$. Since $v_1,\ldots,v_p$ span the domain,
there are scalars $c_j$ with $x=\sum_{j=1}^pc_jv_j$.
Repeated use of linearity gives
\[
 T(x)=\sum_{j=1}^pc_jT(v_j)=\sum_{j=1}^pc_j0=0.
\]
The choice of $x$ was arbitrary, so $T$ sends every input to zero.
The spanning assumption is essential: vanishing on just one line would
not generally determine a map on the whole domain.

\exercise{1.8.31}{A dependence relation survives applying a linear map}
Dependence gives scalars $c_1,c_2,c_3$, not all zero, such that
$c_1v_1+c_2v_2+c_3v_3=0$. Apply $T$ and use $T(0)=0$:
\[
 c_1T(v_1)+c_2T(v_2)+c_3T(v_3)
 =T(c_1v_1+c_2v_2+c_3v_3)=0.
\]
The same coefficients are still not all zero. Thus they form a
nontrivial relation among the three image vectors, proving dependence.

\exercise{1.8.32}{An absolute value violates homogeneity}
For $T(x_1,x_2)=(x_1-2|x_2|,x_1-4x_2)$, choose $u=(0,1)$.
Then $T(u)=(-2,-4)$, so $-T(u)=(2,4)$. But
$T(-u)=T(0,-1)=(-2,4)$, which differs in its first coordinate.
Hence $T(-u)\ne-T(u)$, violating homogeneity for scalar $-1$.
\answer{The transformation is not linear.}

\exercise{1.8.33}{A constant term violates preservation of zero}
The displayed rule gives
$T(0,0)=(0,-3,0)\ne(0,0,0)$.
Since every linear transformation must map the zero input to the zero
output, this single counterexample disproves linearity. No complicated
comparison of general inputs is needed.

\exercise{1.8.34}{Reflection through a coordinate plane is linear}
Let $u,v\in\R^3$ and $c,d\in\R$. For the stated rule,
\begin{align*}
 T(cu+dv)
 &=(cu_1+dv_1,cu_2+dv_2,-cu_3-dv_3)\\
 &=c(u_1,u_2,-u_3)+d(v_1,v_2,-v_3)\\
 &=cT(u)+dT(v).
\end{align*}
This is the combined linearity condition for arbitrary inputs and scalars.
Equivalently, $T$ is multiplication by $\operatorname{diag}(1,1,-1)$.

\exercise{1.8.35}{A coordinate projection is linear}
For $T(x_1,x_2,x_3)=(x_1,0,x_3)$,
\begin{align*}
 T(cu+dv)&=(cu_1+dv_1,0,cu_3+dv_3)\\
 &=c(u_1,0,u_3)+d(v_1,0,v_3)=cT(u)+dT(v).
\end{align*}
The middle zero is preserved under addition and scaling. Thus the map
is linear; its matrix is $\operatorname{diag}(1,0,1)$.

\exercise{1.8.36}{Dependence created in the output forces a nonzero zero-input}
Because $T(u),T(v)$ are dependent, some $a,b$, not both zero,
satisfy $aT(u)+bT(v)=0$. Set $x=au+bv$. Linearity gives
$T(x)=0$. Moreover $x\ne0$: if $au+bv=0$, independence of
$u,v$ would force $a=b=0$, contrary to their choice.
\answer{$x$ is a nontrivial solution of $T(x)=0$.}


\exercise{1.8.37}{Compute all inputs in the kernel}

The supplied matrix is
\[ A=\begin{bmatrix}2 & 3 & 5 & -5\\-7 & 7 & 0 & 0\\-3 & 4 & 1 & 3\\-9 & 3 & -6 & -4\end{bmatrix}.\]

Exact row reduction yields
\[ \operatorname{rref}(A)=\begin{bmatrix}1 & 0 & 1 & 0\\0 & 1 & 1 & 0\\0 & 0 & 0 & 1\\0 & 0 & 0 & 0\end{bmatrix}.\]

Choose the sole free variable to be $t$. Reading the pivot equations gives
\[ x=t\begin{bmatrix}-1\\-1\\1\\0\end{bmatrix},\qquad t\in\R.\]

Direct multiplication of the original $A$ by this generator gives zero in every row. Conversely the pivot equations determine all remaining coordinates from $t$, so the displayed family is the complete solution set.


\exercise{1.8.38}{Compute all inputs in the kernel}

The supplied matrix is
\[ A=\begin{bmatrix}3 & 4 & -7 & 0\\5 & -8 & 7 & 4\\6 & -8 & 6 & 4\\9 & -7 & -2 & 0\end{bmatrix}.\]

Exact row reduction yields
\[ \operatorname{rref}(A)=\begin{bmatrix}1 & 0 & 0 & 1\\0 & 1 & 0 & 1\\0 & 0 & 1 & 1\\0 & 0 & 0 & 0\end{bmatrix}.\]

Choose the sole free variable to be $t$. Reading the pivot equations gives
\[ x=t\begin{bmatrix}-1\\-1\\-1\\1\end{bmatrix},\qquad t\in\R.\]

Direct multiplication of the original $A$ by this generator gives zero in every row. Conversely the pivot equations determine all remaining coordinates from $t$, so the displayed family is the complete solution set.


\exercise{1.8.39}{Find an input for a specified output}

Use the matrix in Exercise 37 and append the stated target. The exact reduced matrix is
\[ [A\mid b]\sim \begin{bmatrix}1 & 0 & 1 & 0 & 1\\0 & 1 & 1 & 0 & 2\\0 & 0 & 0 & 1 & 0\\0 & 0 & 0 & 0 & 0\end{bmatrix}.\]

No row says $0=1$, so the equations are consistent. A particular input is $x=\begin{bmatrix}1\\2\\0\\0\end{bmatrix}$. Substitution gives $Ax=\begin{bmatrix}8\\7\\5\\-3\end{bmatrix}=b$.

All inputs with that same image are
\[ x=\begin{bmatrix}1\\2\\0\\0\end{bmatrix}+t\begin{bmatrix}-1\\-1\\1\\0\end{bmatrix},\quad t\in\R.\]

The homogeneous part changes the input without changing its image, because its product with $A$ is zero.


\exercise{1.8.40}{Find an input for a specified output}

Use the matrix in Exercise 38 and append the stated target. The exact reduced matrix is
\[ [A\mid b]\sim \begin{bmatrix}1 & 0 & 0 & 1 & 1\\0 & 1 & 0 & 1 & 2\\0 & 0 & 1 & 1 & 1\\0 & 0 & 0 & 0 & 0\end{bmatrix}.\]

No row says $0=1$, so the equations are consistent. A particular input is $x=\begin{bmatrix}1\\2\\1\\0\end{bmatrix}$. Substitution gives $Ax=\begin{bmatrix}4\\-4\\-4\\-7\end{bmatrix}=b$.

All inputs with that same image are
\[ x=\begin{bmatrix}1\\2\\1\\0\end{bmatrix}+t\begin{bmatrix}-1\\-1\\-1\\1\end{bmatrix},\quad t\in\R.\]

The homogeneous part changes the input without changing its image, because its product with $A$ is zero.
